% =============================================================================
%  FRONT MATTER: abstract, how to read this book, notation, provenance rules
% =============================================================================

\section*{Abstract}
\label{sec:abstract}
\addcontentsline{toc}{section}{Abstract}

This document is a self-contained, book-length account of a multi-asset statistical
arbitrage project. It is written for a reader who starts from nothing: Part~I develops
every piece of mathematics the project uses --- probability, returns, stationarity,
least squares, regularization, unit-root and cointegration tests, the
Ornstein--Uhlenbeck process, the Kalman filter, nonlinear forecasting, the arithmetic
of backtesting, data-snooping corrections, and the cache/latency arithmetic of a
low-latency C++ engine --- from definitions and first principles, with derivations
rather than citations. Part~II then walks, phase by phase, through what was actually
built and measured in this repository, tying every number in the text to a script that
regenerates it. Part~III synthesises the findings.

The empirical headline is deliberately unglamorous and, we argue, honest. On a real
195-name, 2010--2026 total-return panel of U.S.\ large-cap equities, a causal
walk-forward basket spread does revert: the traded (rolling-hedge) spread rejects a
unit root decisively (ADF $p<10^{-3}$) with an Ornstein--Uhlenbeck half-life of
$\approx 9.4$ trading days, while a static sixteen-year hedge does not
($p\approx0.08$, half-life $\sim 170$--$178$ days). The out-of-sample information
coefficient of the standardized spread against the realized 21-bar reversion return is
positive for every estimator tested ($\approx 0.12$--$0.21$). Ridge shrinkage and,
better, a low-noise Kalman hedge raise the full-sample net Sharpe from
$\approx 0.05$ (rolling OLS) to $\approx 0.09$--$0.11$ at 10\,bps one-way cost, mostly
by cutting turnover rather than by finding more signal. Flexible nonlinear models
(gradient boosting, a shallow network) never beat the regularized linear map
out-of-sample, and a Ramsey RESET test does not reject linearity on the main basket
($p\approx0.38$).

But: after transaction costs the daily edge is small relative to its own noise, the
four-strategy book is net-negative over the full sample (net Sharpe
$\approx -0.40$ to $-0.46$), and --- decisively --- \emph{no} variant's net Sharpe
survives a fair correction for the breadth of the search
($M=44$ audited strategies; Bonferroni survivors $=0$; White Reality Check and
Hansen's SPA both non-significant; every pairwise Diebold--Mariano confidence interval
spans zero). The correct conclusion is that the \emph{structure} is real and
reproducible (mean reversion of an adaptively estimated spread, smooth low-turnover
hedges beating churny refits, linear beating nonlinear net of cost) while the
\emph{magnitude} is not distinguishable from a data-snooping artifact at this sample
size, this frequency, and this cost level. Everything needed to check that claim ---
data pipeline, statistical toolkit with 35 unit tests, a dependency-free C++ engine
with 24 unit tests and 2{,}852 assertions, and the report generator --- is in this
repository and regenerates end to end.

\vspace{1.5em}

\section*{How to read this book}
\label{sec:howtoread}
\addcontentsline{toc}{section}{How to read this book}

The book has three parts with different jobs.

\begin{description}[leftmargin=0pt,style=nextline]
  \item[Part I --- The Mathematics, From Nothing.] Eleven chapters of theory. Nothing
    is assumed beyond high-school algebra and the idea that a probability is a number
    between 0 and 1. Every estimator used later is \emph{derived}: where the formula
    comes from, what assumptions make it valid, and exactly how those assumptions fail
    on financial data. Read it linearly the first time; use it as a reference after
    that. Chapters are ordered so that each only depends on earlier ones:
    \begin{center}
      \begin{tabular}{ll}
        \toprule
        Ch. & Builds on \\
        \midrule
        1 What statistical arbitrage is & --- \\
        2 Probability, returns, stationarity & 1 \\
        3 Linear regression from scratch & 2 \\
        4 Regularization (ridge / lasso / elastic net / PCR) & 3 \\
        5 Time-series econometrics (unit roots, cointegration, HAC) & 2, 3 \\
        6 Mean reversion and the Ornstein--Uhlenbeck process & 2, 5 \\
        7 State-space models and the Kalman filter & 2, 3, 6 \\
        8 Nonlinear models, overfitting, and the RESET test & 3, 4 \\
        9 The mathematics of backtesting (costs, Sharpe, drawdown, bootstrap) & 2, 6 \\
        10 Multiple testing and data snooping & 9 \\
        11 Systems: caches, latency, and incremental linear algebra & 3 \\
        \bottomrule
      \end{tabular}
    \end{center}
  \item[Part II --- The Project, Phase by Phase.] Fourteen chapters: a phase map, then
    one chapter per phase (``Part'') of the build plan in \code{PLAN.md} (Phases 12 and 13,
    both assembly and review, share the last one). Each chapter has the same shape:
    \emph{goal} $\rightarrow$ \emph{what was built} (with the files that implement it)
    $\rightarrow$ \emph{the math it uses} (pointing back to Part~I) $\rightarrow$
    \emph{results} (tables and figures regenerated by code) $\rightarrow$
    \emph{honest reading} $\rightarrow$ \emph{exit gate}. If you want to know ``what did
    you actually do, and in what order?'', this is the part.
  \item[Part III --- Synthesis.] The combined verdict, what this project does
    \emph{not} claim (non-goals), and the concrete next steps that would be needed to
    turn a real-in-structure, too-small-in-magnitude edge into something tradeable.
  \item[Appendices.] Four reference sections: twelve extended derivations that would
    interrupt the flow of Part~I (Appendix~\ref{ch:derivations}); the validation suite and
    all 35 unit tests of the statistical toolkit (Appendix~\ref{ch:statkitappendix}); a
    reproducibility manual with the exact commands, their expected console output, and the
    map from every table in this book to the script that produces it
    (Appendix~\ref{ch:reproducibility}); and a symbol index, glossary, list of figures and
    list of tables (Appendix~\ref{ch:glossary}).
\end{description}

\paragraph{Boxes.} Three kinds of aside appear throughout.
\emph{Intuition} boxes give the plain-English picture before or after the algebra.
\emph{Derivation} material is inline and marked by equation numbers.
\emph{Honesty} remarks flag places where a shortcut, an assumption, or a negative
result is being reported rather than hidden --- they are the most important paragraphs
in the book.

\paragraph{Two formats, one source.} The canonical source is
\code{report/report.tex}, which \code{\textbackslash input}s the 29 chapter files in
\code{report/chapters/} and compiles to a paginated PDF with \code{latexmk -pdf} wherever
a \TeX\ distribution exists. In this sandbox none is installable, so
\code{python3 report/build\_report.py} also renders the same source to
\code{report/report.html}: a linked, navigable book whose mathematics is rasterised locally
with matplotlib's \code{mathtext} (no CDN, no MathJax, no network). The renderer is
\code{report/latex\_to\_html.py}, described in Section~\ref{sec:phase13:renderer}; it is
a preview, not a \TeX\ implementation, and the differences are cosmetic (floats are placed
where they are written, page breaking is ignored). If you read only one format, read the
HTML here and the PDF anywhere else --- the numbers are identical because both come from the
same generated tables.

\paragraph{Provenance rule.} No number appears in this book unless it was produced by
code checked into this repository. Figures and tables are generated into
\code{results/} and \code{report/tables/} by the scripts named in each chapter's
``reproduce it'' line, and the whole report is rebuilt with
\code{python3 report/build_report.py}. Where a number in the prose is quoted for
readability, the generating table is \code{\textbackslash input}-ed on the same page so
you can check it.

\vspace{1em}

\section*{Notation}
\label{sec:notation}
\addcontentsline{toc}{section}{Notation}

\begin{center}
\renewcommand{\arraystretch}{1.25}
\begin{tabular}{lp{0.72\linewidth}}
\toprule
Symbol & Meaning \\
\midrule
$T$ & number of time observations (bars) in a sample \\
$N,\ n$ & number of assets / number of regressors (features) \\
$t$ & time index, $t=1,\dots,T$; bars are daily unless stated \\
$P_{i,t}$ & price of asset $i$ at $t$; $P^{T}$ is the \emph{target}, $P^{j}$ a basket name \\
$p_{i,t}$ & log price, $\log P_{i,t}$ \\
$r_{i,t}$ & simple return, $P_{i,t}/P_{i,t-1}-1$ \\
$y_t,\ x_t$ & regression target and regressor row vector at time $t$ \\
$Y,\ X$ & stacked sample: $Y\in\R^{T}$, $X\in\R^{T\times n}$ \\
$\beta$ & vector of hedge weights (regression coefficients) \\
$\hat\beta_t$ & estimate of $\beta$ using information through bar $t$ only \\
$\eps_t,\ \hat\eps_t$ & model error and estimated residual \\
$S_t$ & spread: $S_t = y_t - \hat\alpha_t - x_t^{\top}\hat\beta_t$ \\
$\mu_S,\ \sigma_S$ & mean and standard deviation of the spread over the estimation window \\
$z_t$ & standardized spread, $z_t=(S_t-\mu_S)/\sigma_S$ \\
$\lambda,\ \alpha$ & regularization penalty strength (context makes clear which) \\
$\theta,\ \mu,\ \sigma$ & OU mean-reversion speed, long-run mean, diffusion \\
$\tau_{1/2}$ & mean-reversion half-life, $\ln 2/\theta$ \\
$Q,\ R$ & Kalman state-noise and observation-noise covariances \\
$\beta_{t\mid t}$ & Kalman posterior (filtered) state at $t$; $\beta_{t\mid t-1}$ the prediction \\
$w_i$ & portfolio weight on strategy or name $i$ \\
$c$ & one-way transaction cost, in basis points of traded notional \\
$g$ & grossing factor, $g = 1 + \sum_i \lvert \beta_i\rvert$ \\
$\mathrm{IC}$ & information coefficient: correlation of a signal with a realized forward return \\
$\hat\Omega$ & long-run (HAC) covariance matrix of a moment condition \\
$M$ & number of hypotheses/strategies searched (the data-snooping count) \\
$\E[\cdot],\ \Var(\cdot),\ \Cov(\cdot)$ & expectation, variance, covariance \\
$\xrightarrow{\ d\ }$ & convergence in distribution \\
$\mathcal{N}(m,\Sigma)$ & Gaussian with mean $m$ and covariance $\Sigma$ \\
\bottomrule
\end{tabular}
\end{center}

Time indices matter more than any other notational choice in this book. A quantity
with a subscript $t$ that is used to make a decision at $t$ must be computable from
data up to and including $t$ --- and, because we execute on the \emph{next} bar, the
hedge that scores bar $t$ is estimated through bar $t-1$. We write this explicitly
wherever it matters, and one of the C++ unit tests exists solely to assert it.
